\subsection{\texorpdfstring{\(\tau\)-Extensions of Groups}{\textbackslash tau-Extensions of Groups}}

In this subsection, we fix a group G, an abelian group F written multiplicatively, and a group homomorphism \(\tau\) from G to the automorphism group \(\operatorname{Aut}(F)\) of F. We denote the identity element of G by e and the identity element of F by 1.

Recall (I, §6, No.~1, p.~65) that an extension \(\mathcal{E}\) of \(G\) by \(F\) is a triple \((\Gamma, \iota, \pi)\) , where \(\Gamma\) is a group, \(\pi: \Gamma \to G\) is a surjective homomorphism, and \(\iota\) is an injective homomorphism from \(F\) to \(\Gamma\) such that \(\operatorname{Im}(\iota) = \operatorname{Ker}(\pi)\) . Let \(\mathcal{E} = (\Gamma, \iota, \pi)\) be such an extension. For every \(\gamma \in \Gamma\) , the mapping \(\varphi_{\gamma}: F \to F\) defined by

\[
\iota (\varphi_ {\gamma} (f)) = \gamma \iota (f) \gamma^ {- 1}
\]

for \(f \in \mathrm{F}\) is an automorphism of \(\mathrm{F}\) . Since \(\mathrm{F}\) is commutative, for every \(f \in \mathrm{F}\) , the automorphism defined by \(\iota(f)\) is the identity mapping on \(\mathrm{F}\) . By passing to the quotient, we obtain a homomorphism \(\operatorname{Int}_{\mathcal{E}}\) from \(\mathrm{G}\) to \(\operatorname{Aut}(\mathrm{F})\) characterized by

\[
\gamma \iota (f) \gamma^ {- 1} = \iota (\mathrm{Int} _ {\mathcal {E}} (\pi (\gamma)) \cdot f)
\]

for \(\gamma \in \Gamma\) and \(f\in \mathrm{F}\)

A \(\tau\) -extension of G by F is an extension \(\mathcal{E} = (\Gamma, \iota, \pi)\) such that \(Int_{\mathcal{E}}\) is equal to \(\tau\) , in other words, that satisfies the relation

\[
\gamma \iota (f) \gamma^ {- 1} = \iota (\tau (\pi (\gamma)) \cdot f)\tag{1}
\]

for \(\gamma \in \Gamma\) and \(f \in \mathrm{F}\) . If \(\mathcal{E} = (\Gamma, \iota, \pi)\) and \(\mathcal{E}' = (\Gamma', \iota', \pi')\) are \(\tau\) -extensions, then a morphism of \(\tau\) -extensions from \(\mathcal{E}\) to \(\mathcal{E}'\) is a morphism of extensions from \(\mathcal{E}\) to \(\mathcal{E}'\) (I, §6, No.~1, p.~65), that is, a group homomorphism \(u: \Gamma \to \Gamma'\) such that \(\pi' \circ u = \pi\) and \(\iota' = u \circ \iota\) . Note that the \(\tau\) -extensions form a species of structures for which every morphism is an isomorphism (Set Theory, IV, §1, No.~5, p.~264 and I, §6, No.~1, p.~66, Proposition 1). Denote by \(\operatorname{Iso}_{\tau} \{\mathcal{E}, \mathcal{F}\}\) the relation

`` \(\mathcal{E}\) and \(\mathcal{F}\) are isomorphic \(\tau\) -extensions.''

This is an equivalence relation; the class of the extension \(\mathcal{E}\) is the class of objects equivalent to \(\mathcal{E}\) for \(\mathrm{Iso}_{\tau}\) (Set Theory, II, §6, No.~6, p.~122).

Lemma 1. --- The relation

\textbf{\texorpdfstring{``α is the class of a τ-extension for \(Iso_{\tau}\) ''}{``α is the class of a τ-extension for Iso\_\{\textbackslash tau\} ''}}

is collectivizing in \(\alpha\) .

Set \(E_{0}=F\times G\) , and consider the mappings \(\iota_{0}:F\to E_{0}\) and \(\pi_{0}:E_{0}\to G\) defined by \(\iota_{0}(f)=(f,e)\) for \(f\in F\) and \(\pi_{0}(f,g)=g\) for \((f,g)\in F\times G\) . Let \(\mathcal{E}=(\Gamma,\iota,\pi)\) be a \(\tau\) -extension of G by F. The mapping \(\pi\) is surjective, so has a section \(\sigma:G\to\Gamma\) such that \(\sigma(e)=e\) . The mapping \(u:(f,g)\mapsto\iota(f)\sigma(g)\) from \(F\times G\) to \(\Gamma\) is bijective. We endow \(F\times G\) with the group law obtained by transfer of structure. The triple \((\mathrm{E}_{0},\iota_{0},\pi_{0})\) is then a \(\tau\) -extension isomorphic to E. Lemma 1 now follows from Set Theory, II, §6, No.~6, p.~122.

We denote by \(\operatorname{Ex}_{\tau}(\mathbf{G},\mathbf{F})\) the set of classes of \(\tau\) -extensions for the relation \(Iso_{\tau}\) .

Examples. --- 1) The external semidirect product \((\mathrm{F} \times_{\tau} \mathrm{G}, i, p)\) (I, §6, No.~1, p.~66, Proposition 3) is a \(\tau\) -extension; we denote it by \(\mathcal{I}_{\tau}\) . We say that a \(\tau\) -extension is semitrivial if it is isomorphic to the \(\tau\) -extension \(\mathcal{I}_{\tau}\) .

\begin{enumerate}
\def\labelenumi{\arabic{enumi})}
\setcounter{enumi}{1}
\tightlist
\item
  For \(i \in \{1,2\}\) , take a group \(G_{i}\) , an abelian group \(F_{i}\) , and a group homomorphism \(\tau_{i}\) from \(G_{i}\) to the automorphism group of \(F_{i}\) . We denote by \(\tau_{1} \times \tau_{2}: G_{1} \times G_{2} \to \text{Aut}(F_{1} \times F_{2})\) the homomorphism defined by
\end{enumerate}

\[
(\tau_ {1} \times \tau_ {2}) (g _ {1}, g _ {2}) \cdot (f _ {1}, f _ {2}) = (\tau_ {1} (g _ {1}) \cdot f _ {1}, \tau_ {2} (g _ {2}) \cdot f _ {2})
\]

for \(g_{1} \in \mathrm{G}_{1}\) , \(g_{2} \in \mathrm{G}_{2}\) , \(f_{1} \in \mathrm{F}_{1}\) , and \(f_{2} \in \mathrm{F}_{2}\) . Let \(\mathcal{E}_1 = (\Gamma_1, \iota_1, \pi_1)\) be a \(\tau_1\) -extension of \(\mathrm{G}_1\) by \(\mathrm{F}_1\) and \(\mathcal{E}_2 = (\Gamma_2, \iota_2, \pi_2)\) a \(\tau_2\) -extension of \(\mathrm{G}_2\) by \(\mathrm{F}_2\) . Then the triple \((\Gamma_1 \times \Gamma_2, \iota_1 \times \iota_2, \pi_1 \times \pi_2)\) is a \(\tau_1 \times \tau_2\) -extension of \(\mathrm{G}_1 \times \mathrm{G}_2\) by \(\mathrm{F}_1 \times \mathrm{F}_2\) ; it is called the exterior product of the extensions \(\mathcal{E}_1\) and \(\mathcal{E}_2\) and denoted by \(\mathcal{E}_1 \times \mathcal{E}_2\) .

